% !TEX program = tectonic
\documentclass[9pt,a4paper]{extarticle}
\usepackage[top=0.8cm,bottom=0.8cm,left=0.9cm,right=0.9cm]{geometry}
\usepackage[T1]{fontenc}\usepackage[utf8]{inputenc}\usepackage{lmodern}
\usepackage[french]{babel}
\usepackage{amsmath,amssymb}\usepackage{textcomp}\usepackage{xcolor}\usepackage{booktabs}\usepackage{array}
\newcolumntype{R}[1]{>{\raggedright\arraybackslash}p{#1}}
\usepackage{enumitem}\usepackage{mdframed}\usepackage{multicol}\usepackage{graphicx}
\graphicspath{{figs_nb11/}}
\definecolor{bleu}{RGB}{20,77,144}\definecolor{vert}{RGB}{20,110,60}
\definecolor{rouge}{RGB}{160,30,30}\definecolor{grisf}{RGB}{110,110,110}
\newcommand{\fig}[2]{\begin{center}\includegraphics[width=\linewidth]{#1}\end{center}\vspace{-2.5mm}{\scriptsize\color{grisf} #2}\par\vspace{0.3mm}}
\pagestyle{empty}\setlength{\parindent}{0pt}\setlength{\parskip}{1.1pt}
\setlength{\abovedisplayskip}{2.5pt}\setlength{\belowdisplayskip}{2.5pt}
\setlength{\abovedisplayshortskip}{1pt}\setlength{\belowdisplayshortskip}{1pt}
\renewcommand{\arraystretch}{0.95}
\setlength{\columnsep}{6mm}\setlength{\columnseprule}{0.3pt}
\renewcommand{\columnseprulecolor}{\color{grisf!40}}
\newcommand{\et}[2]{\vspace{1.3mm}{\color{bleu}\bfseries #1 $\cdot$ #2}\par\vspace{-0.2mm}}
\newcommand{\res}[1]{\vspace{0.3mm}{\small\color{bleu}$\Rightarrow$\ #1}\par\vspace{0.2mm}}
\newcommand{\obs}[1]{\vspace{0.2mm}\hangindent=3.4mm\hangafter=0{\color{bleu}\footnotesize\textbullet}~{\small #1}\par\vspace{0.2mm}}
\newcommand{\lire}[1]{{\scriptsize\color{vert}$\blacktriangleright$~\textbf{lire :} #1}\par}
\newenvironment{pts}{\begin{itemize}[leftmargin=3.4mm,itemsep=0.2mm,topsep=0.4mm,parsep=0pt,label={\color{bleu}\footnotesize\textbullet}]}{\end{itemize}}
\newmdenv[linecolor=bleu!55,backgroundcolor=bleu!5,linewidth=0.5pt,innertopmargin=2.5pt,
 innerbottommargin=2.5pt,innerleftmargin=3pt,innerrightmargin=3pt,skipabove=1.5pt,skipbelow=1.5pt]{boite}
\begin{document}

\begin{center}
{\large\bfseries Stratégie par ticker --- suite de tests}
\end{center}
\vspace{-1.5mm}

\begin{multicols}{2}

% ═══════════ É1 ═══════════
\et{É1}{La donnée --- le périmètre}
\begin{pts}
\item On part de la table déjà propre (pipeline S1/S2) --- aucun nettoyage refait.
\item On ne garde qu'un sous-ensemble : actions cotées, 2014--2026, série de prix exploitable.
\end{pts}
\begin{center}\scriptsize
\begin{tabular}{@{}l r l@{}}
\toprule
table clean & 134\,464 & \\
$\to$ actions seules & 128\,974 & {\color{grisf}$-$5\,490 : 2\,190 ETF, 3\,300 inconnu} \\
$\to$ années 2014--2026 & 128\,269 & {\color{grisf}$-$705 : donnée trop mince} \\
$\to$ prix exploitable & \textbf{113\,369} & {\color{grisf}$-$14\,900} \\
\bottomrule
\end{tabular}
\end{center}
\lire{chaque ligne retire un morceau ; la colonne = nombre de transactions qui restent.}
{\scriptsize\color{grisf} 2\,755 titres $\cdot$ 350 membres $\cdot$ 57\,445 achats, 55\,924 ventes.\par}

% ═══════════ É2 ═══════════
\et{É2}{Dater les lots (FIFO) et purger le dormant}
\begin{pts}
\item Chaque achat empile un lot : (date, nombre de parts).
\item Une vente consomme les lots \textbf{les plus anciens} d'abord.
\item L'âge ne choisit pas le lot vendu, il l'\textbf{étiquette} : \og récent\fg{} si détenu $\le 2$ ans, \og vieux\fg{} sinon.
\item Le signal ne garde que la part \textbf{récente} d'une vente (le dormant ne signale pas un délit d'initié).
\end{pts}
\[ \gamma_k = \frac{\sum_{\ell \subset k} \mathbf{1}\{h_\ell \le \bar a\}\, q_\ell}{\sum_{\ell \subset k} q_\ell} \in [0,1] \]
\begin{pts}
\item $s_k$ : jour de bourse où le trade $k$ s'exécute.
\item $q_k = A_k / P(s_k)$ : nombre de parts ($A_k$ = montant, $P$ = prix).
\item $h_\ell$ : durée de détention du lot $\ell$ vendu, en jours de bourse.
\item $\bar a = 504$ j : seuil \og vieux\fg{} = 2 ans de bourse.
\item $\gamma_k$ : part de la vente $k$ portant sur du récent ($0$ = tout vieux, $1$ = tout récent).
\item $\mathbf{1}\{\cdot\}$ = 1 si la condition est vraie, 0 sinon $\cdot$ $\ell \subset k$ : lot $\ell$ (de $q_\ell$ parts) consommé par la vente $k$.
\end{pts}
\fig{durees.png}{}
\lire{$x$ = durée de détention des ventes (jours) ; barre rouge = $\bar a$ (2 ans) ; à droite = lots dormants, retirés du signal.}
\begin{center}\scriptsize
\begin{tabular}{@{}l r@{}}
\toprule
$\gamma{=}1$ vente entièrement récente (gardée) & 79,9 \% \\
$\gamma{=}0$ entièrement vieille (retirée) & \textbf{17,1 \%} \\
$0{<}\gamma{<}1$ partielle & 3,1 \% \\
\midrule
masse vendue gardée dans le signal & \textbf{89,5 \%} des 2,11 Md\$ \\
\bottomrule
\end{tabular}
\end{center}
\lire{les 3 premières lignes = quelle fraction des ventes tombe dans ce cas ; dernière = part du \$ vendu qui alimente le signal.}

% ═══════════ É3 ═══════════
\et{É3}{Le signal ticker et la pondération}
\begin{pts}
\item Sur une fenêtre de $W$ jours, on somme les achats et la part récente des ventes de chaque titre.
\item On retient un titre s'il est net-acheté par \textbf{au moins deux membres}.
\item Chaque membre pèse pareil ($1/K_t$), qu'il ait mis 11 M\$ ou 15 k\$, et répartit son poids entre ses titres.
\end{pts}
\[ NetBuy_{i,t} = Buy_{i,t} - Sell_{i,t}\cdot\gamma \]
\[ S_t = \{\, i : NetBuy_{i,t} > 0 \;\wedge\; nBuyers_{i,t} \ge 2 \,\} \]
\[ w_{i,t} = \frac{1}{K_t} \sum_{m \in \mathcal{M}_t} \frac{\mathbf{1}\{ i \in B_{m,t} \}}{|B_{m,t}|} \]
\begin{pts}
\item $Buy,\ Sell$ : montants achetés / vendus du titre $i$ dans la fenêtre.
\item $nBuyers$ : nombre de membres distincts ayant acheté $i$.
\item $S_t$ : les titres retenus le mois $t$ ; $B_{m,t}$ : ceux de $S_t$ achetés par le membre $m$.
\item $\mathcal{M}_t$ : les membres ayant acheté $\ge1$ titre de $S_t$ (membres actifs) ; $K_t = |\mathcal{M}_t|$ ; $w_{i,t}$ : poids du titre $i$ ($\sum_i w_i = 1$).
\end{pts}

% ═══════════ É4 ═══════════
\et{É4}{Choisir la fenêtre $W$ (test de concentration)}
\begin{pts}
\item On teste 4 fenêtres et on tranche sur la \textbf{concentration}, jamais sur le rendement.
\end{pts}
\[ W^\star = \min\bigl\{\, W : \mathbb{P}[\max_i w_{i,t} > 10\,\%] \le 20\,\% \,\bigr\} \]
\begin{pts}
\item $\max_i w_{i,t}$ : poids de la plus grosse ligne, un mois donné (règle : $>10\%$ dans $\le20\%$ des mois).
\end{pts}
\begin{center}\scriptsize
\begin{tabular}{@{}l r r r r@{}}
\toprule
$W$ & 21 j & \textbf{42 j} & 63 j & 126 j \\
\midrule
$|S_t|$ médian & 28 & \textbf{64} & 91 & 156 \\
$K_t$ médian & 19 & \textbf{32} & 39 & 53 \\
$\max_i w_i$ médian & 11,0 \% & \textbf{7,2 \%} & 6,3 \% & 5,4 \% \\
mois $> 10\,\%$ & 58 \% & \textbf{17 \%} & 9 \% & 5 \% \\
\bottomrule
\end{tabular}
\end{center}
\lire{colonnes = les 4 fenêtres testées ; $|S_t|$ = titres retenus/mois ; $K_t$ = membres actifs/mois ; $\max w$ = plus grosse ligne ; dernière ligne = fréquence de sur-concentration (le critère).}
\res{$W = 42$ j retenu.}
\obs{à $W{=}21$ une ligne montait à 50 \% et $K_t$ tombait à 2 ; à $W{=}42$, $|S_t|$ médian 64, $K_t$ 32, jamais vide.}

% ═══════════ encadré lecture ═══════════
\begin{boite}
{\footnotesize\bfseries\color{bleu}Comment lire les résultats} {\scriptsize\color{grisf}--- notations.}
\begin{pts}
\item \textbf{excès/an} : $\displaystyle x_Y = \prod_{t\in Y}(1{+}r_t) - \prod_{t\in Y}(1{+}r_{\mathrm{SPY},t})$ pour l'année $Y$, puis moyenne sur les années.
\item \textbf{$t$} : $t = \bar x \,/\, (\sigma_x/\sqrt{n})$, $n$ = nombre d'années ($|t|\ge2$ : significatif à 5 \%).
\item \textbf{$\alpha,\ \beta$} : régression $r_t - r_{f,t} = \alpha + \beta\,(r_{\mathrm{SPY},t} - r_{f,t}) + \varepsilon_t$, annualisé $\alpha_{\mathrm{an}} = 252\,\alpha$ ($r_f$ = taux sans risque, $\varepsilon$ = résidu).
\item \textbf{NAV} : $\mathrm{NAV}_T = 100\prod_{t\le T}(1{+}r_t)$, base 100 en 2014 (SPY : 517).
\end{pts}
\end{boite}

% ═══════════ É5 ═══════════
\et{É5}{Test 1 --- la base}
\begin{pts}
\item Long seul (on achète, jamais de vente à découvert).
\item Fenêtre du signal sur la \textbf{transaction} $\tau$ (date où le membre a agi).
\item Aucune contrainte de rotation. Coupe en fin de mois, exécution au \textbf{close J+1} (cours de clôture du lendemain).
\end{pts}
\res{excès $\mathbf{+1{,}71\,\%}$/an $\cdot$ $t\,1{,}03$ $\cdot$ $\alpha\,+1{,}05\,\%$ $\cdot$ $\beta\,1{,}02$ $\cdot$ NAV \textbf{622}}
\obs{7/13 années positives ; $\beta \approx 1$ ; NAV 622 $>$ SPY 517.}
\fig{beta_glissant.png}{}
\lire{courbe bleue = $\beta$ de la base sur 1 an glissant ; il reste entre \textbf{0,93 et 1,11} (médiane 1,02) ; pointillés = $\beta{=}1$ (marché) et $\beta$ plein-échantillon 1,02.}

% ═══════════ É6 ═══════════
\et{É6}{Test 2 --- + cap de rotation 25 \%/mois}
\begin{pts}
\item On ne réalloue pas d'un coup : on avance \emph{partiellement} vers la cible, d'au plus 25 \%/mois.
\end{pts}
\[ \mathrm{TO} = \tfrac12 \textstyle\sum_i |w^{\text{cible}}_i - w^{\text{avant}}_i| \]
\[ \lambda = \min\!\bigl(1,\ \tau_{\max}/\mathrm{TO}\bigr), \qquad \tau_{\max}{=}25\,\% \]
\begin{pts}
\item $\mathrm{TO}$ : rotation voulue ce mois ($0$ = rien ne bouge, $1$ = tout change).
\item $\lambda$ : fraction du chemin vers la cible qu'on parcourt réellement.
\end{pts}
\res{excès $\mathbf{+0{,}27\,\%}$/an $\cdot$ $t\,0{,}23$ $\cdot$ $\alpha\,-0{,}16\,\%$ $\cdot$ $\beta\,1{,}01$ $\cdot$ NAV \textbf{531}}
\obs{le cap ramène la rotation de 59 \% à 25 \%/mois ; NAV 531 (base 622) ; excès $+0{,}27\,\%$/an ($t\,0{,}23$) contre $+1{,}71\,\%$ ($t\,1{,}03$).}

% ═══════════ É7 ═══════════
\et{É7}{Test 3 --- + jambe short (long-short)}
\begin{pts}
\item Au lieu d'ignorer les titres net-\textbf{vendus}, on les shorte (on parie sur leur baisse).
\item On ne short que ceux corroborés par $\ge2$ vendeurs de stock récent.
\end{pts}
\[ r = r^{+} - \theta\, r^{-}, \qquad \theta = 1 \]
\begin{pts}
\item $S^{-}$ : titres à flux net négatif, $\ge2$ vendeurs récents ; $r^{+},r^{-}$ : rendement du livre long / short.
\item $\theta$ : intensité du short ($\theta{=}0$ redonne le long seul).
\end{pts}
\res{excès $\mathbf{-12{,}8\,\%}$/an $\cdot$ $t\,-3{,}19$ $\cdot$ $\beta\,-0{,}02$ $\cdot$ NAV \textbf{122}}
\obs{$|S^{-}|$ médian 45 ; les titres net-vendus \emph{montent} après la vente ; $\beta$ passe à $\sim$0.}
\fig{extensions.png}{}
\lire{$y$ = NAV base 100 (échelle log) ; A = base (T1), B = $+$cap 25 \% (colle à A), C = long-short (s'effondre), gris = SPY.}

% ═══════════ É8 short-only ═══════════
\et{É8}{Test 4 --- short-only (la jambe short seule)}
\begin{pts}
\item On enlève toute position longue : on \textbf{shorte} seulement les titres net-vendus ($\ge2$ vendeurs récents).
\end{pts}
\res{excès $\mathbf{-26{,}6\,\%}$/an $\cdot$ $t\,-3{,}40$ $\cdot$ $\alpha\,-2{,}36\,\%$ $\cdot$ $\beta\,\mathbf{-1{,}03}$ $\cdot$ NAV \textbf{16}}
\obs{les titres net-vendus montent : shortés seuls, 100~\$ tombent à 16 ; $\beta\,{\approx}\,{-}1$ = à l'inverse du marché.}

% ═══════════ É9 δ ═══════════
\et{É9}{Test 5 --- à la date de divulgation $\delta$}
\begin{pts}
\item Même base (long seul), mais on agit à la \textbf{divulgation} $\delta$ (quand le trade devient public) au lieu de $\tau$.
\item Le FIFO et $\gamma$ restent sur $\tau$ : la durée de détention est une réalité, pas une date de publication.
\end{pts}
\res{excès $\mathbf{+0{,}35\,\%}$/an $\cdot$ $t\,0{,}18$ $\cdot$ $\alpha\,-0{,}03\,\%$ $\cdot$ $\beta\,1{,}01$ $\cdot$ NAV \textbf{530}}
\obs{lag médian de divulgation 27 j ; l'écart cumulé sur SPY passe de $+105$ pts ($\tau$) à $+13$ pts ($\delta$).}
\begin{center}\scriptsize
\begin{tabular}{@{}l r r@{}}
\toprule
excès/an (NAV) & sans cap & cap 25 \% \\
\midrule
$\tau$ (transaction) & $+1{,}71\,\%$ (622) & $+0{,}27\,\%$ (531) \\
$\delta$ (divulgation) & $+0{,}35\,\%$ (530) & $-0{,}49\,\%$ (487) \\
\bottomrule
\end{tabular}
\end{center}
{\color{bleu}\footnotesize\textbullet}~{\small le cap réduit l'excès à $\tau$ comme à $\delta$.\par}
\fig{coins.png}{}
\lire{$y$ = NAV base 100 (log) ; $\tau$ = transaction, $\delta$ = divulgation ; trait plein = sans cap, tireté = cap 25 \% ; pointillé = SPY.}
\begin{center}\scriptsize
\begin{tabular}{@{}l r r@{\quad}l r r@{}}
\toprule
année & exc.\,$\tau$ & exc.\,$\delta$ & année & exc.\,$\tau$ & exc.\,$\delta$ \\
\midrule
2014 & $-1{,}3$ & $+2{,}2$ & 2021 & $+2{,}0$ & $\mathbf{+13{,}2}$ \\
2015 & $-0{,}9$ & $+0{,}1$ & 2022 & $-3{,}1$ & $-1{,}6$ \\
2016 & $+9{,}4$ & $+11{,}1$ & 2023 & $\mathbf{+11{,}5}$ & $-3{,}8$ \\
2017 & $+3{,}0$ & $+3{,}7$ & 2024 & $-8{,}7$ & $-9{,}4$ \\
2018 & $+0{,}9$ & $-2{,}4$ & 2025 & $+10{,}9$ & $+3{,}5$ \\
2019 & $-3{,}4$ & $+3{,}0$ & 2026 & $-1{,}6$ & $-5{,}4$ \\
2020 & $+3{,}5$ & $\mathbf{-9{,}6}$ & \textbf{moy.} & $\mathbf{+1{,}7}$ & $\mathbf{+0{,}4}$ \\
\bottomrule
\end{tabular}
\end{center}
\lire{une ligne = une année ; exc.\,$\tau$ et exc.\,$\delta$ = l'excès sur SPY selon qu'on agit à la transaction ou à la divulgation ; moy. = la colonne excès/an du récap.}

% ═══════════ récap ═══════════
\et{}{Les cinq tests côte à côte}
\begin{center}\scriptsize
\begin{tabular}{@{}l r r r r r@{}}
\toprule
test & excès/an & $t$ & $\alpha$/an & $\beta$ & NAV \\
\midrule
T1 base (long, $\tau$) & $+1{,}71\,\%$ & 1,03 & $+1{,}05\,\%$ & 1,02 & 622 \\
T2 $+$ cap 25 \% & $+0{,}27\,\%$ & 0,23 & $-0{,}16\,\%$ & 1,01 & 531 \\
T3 $+$ long-short ($\theta{=}1$) & $-12{,}8\,\%$ & $-3{,}19$ & $-0{,}25\,\%$ & $-0{,}02$ & 122 \\
T4 short-only ($\tau$) & $-26{,}6\,\%$ & $-3{,}40$ & $-2{,}36\,\%$ & $-1{,}03$ & 16 \\
T5 base, à $\delta$ & $+0{,}35\,\%$ & 0,18 & $-0{,}03\,\%$ & 1,01 & 530 \\
\bottomrule
\end{tabular}
\end{center}
\lire{les 5 colonnes = cf. encadré \og Comment lire les résultats\fg{}.}

% ─── É10 · Test 6 · Comités ───
\et{É10}{Test 6 --- comités pertinents}
\textbf{Ce qu'on fait} (le reste est la base à l'identique, \textbf{en long}) :
\begin{pts}
\item \textbf{Le filtre} : on ne garde que les achats d'un membre \textbf{siégeant sur le comité qui régule le secteur du titre} --- Financial Services$\to$Financials, Armed Services$\to$Industriels, Energy \& Commerce$\to$\{Énergie, Santé\}$\dots$ ($18\,948$ trades, $16{,}7\,\%$).
\item On \textbf{achète} $S^{c}_t$ (aucune vente à découvert), une voix par membre restreinte à ces membres, mensuel, close J+1. \emph{Variante grossière} : comité \og clé\fg{} (drapeau, sans le secteur).
\end{pts}
\[ k \text{ retenu} \iff \text{secteur}(i_k) \in \text{secteurs}\bigl(\text{comités}(m_k)\bigr) \]
\[ w^{c}_{i,t} = \frac{1}{K^{c}_t}\sum_{m\in\mathcal{M}^{c}_t}\frac{\mathbf{1}\{i\in B^{c}_{m,t}\}}{|B^{c}_{m,t}|}, \qquad \textstyle\sum_i w^{c}_i = 1 \]
\begin{center}\scriptsize
\begin{tabular}{@{}l r r r r@{}}
\toprule
& exc.\,$\tau$ & $t_\tau$ & exc.\,$\delta$ & $t_\delta$ \\
\midrule
comité-secteur & $\mathbf{+7{,}67}$ & \textbf{1,60} & $-3{,}65$ & $-0{,}99$ \\
comité-clé (drapeau) & $+1{,}74$ & 0,70 & $-1{,}05$ & $-0{,}56$ \\
\bottomrule
\end{tabular}
\end{center}
\fig{comites.png}{}
\textbf{Ce qu'on observe} :
\begin{pts}
\item À la transaction $\tau$, restreindre aux comités \textbf{double la base} : NAV \textbf{1102} contre 622, $\alpha\,{+}5{,}95\,\%$/an ($t_\alpha$ 1,35).
\item Il \textbf{s'effondre à la divulgation} $\delta$ : NAV 326, $-3{,}65\,\%$/an.
\item \textbf{À nuancer} : seulement $\sim$4 titres/mois --- l'excès annuel part dans tous les sens, et le filtre est très déséquilibré par secteur (ci-dessous).
\end{pts}
\fig{comites_annuel.png}{excès annuel à $\tau$ (bleu) et $\delta$ (vert) : de $-30$ à $+45\,\%$ selon l'année.}
\fig{comites_secteurs.png}{\% des trades de chaque secteur retenus par le filtre : 63 \% Industrie, 0 \% Cons.\,Disc./Immobilier.}

% ═══════════ Produits réels : NANC / GOP / Quiver ═══════════
\et{}{Les produits réels --- ce qu'ils annoncent vs ce qu'on obtient}
Trois produits commercialisent la \og{}copie du Congrès\fg{}. Pour chacun : la règle qu'ils \textbf{disent} suivre, le chiffre qu'ils \textbf{annoncent}, puis \textbf{notre} réplique fidèle de leur règle.

\textbf{1. Quiver Quantitative} (indices \og{}Congress Trading\fg{}). \textbf{Ce qu'ils disent faire} :
\begin{pts}
\item \textbf{Congress Buys} : long les titres \textbf{achetés} par le Congrès, \textbf{pondérés à la taille du trade} (\$), rotation \textbf{hebdo}. \textbf{Long-Short 130/30} : $+$130 \% des achats, $-$30 \% des ventes.
\item On réplique à la transaction $\tau$ (borne haute) \emph{et} à la divulgation $\delta$ (réplicable) ; fenêtre 63 j, rotation mensuelle.
\end{pts}
\[ w^{\text{Buys}}_{i,t} \;=\; \frac{\displaystyle\sum_{k\,:\ achat,\ i_k=i,\ t-W<s_k\le t} A_k}{\displaystyle\sum_{j}\ \sum_{k\,:\ achat,\ i_k=j,\ t-W<s_k\le t} A_k},\qquad \textstyle\sum_i w^{\text{Buys}}_i=1 \]
\[ r^{\text{LS}}_t \;=\; 1{,}3\,r^{\text{long}}_t - 0{,}3\,r^{\text{short}}_t \]
\begin{center}\scriptsize
\begin{tabular}{@{}l r r r@{}}
\toprule
CAGR/an \emph{2020-26} & annoncé & rép.\,$\tau$ & rép.\,$\delta$ \\
\midrule
Congress Buys & 36,2 \% \, $\beta$1,14 & \textbf{22,6 \%} & 21,4 \% \\
Long-Short 130/30 & 34,0 \% \, $\beta$1,15 & \textbf{23,5 \%} & 22,3 \% \\
SPY (même fenêtre) & --- & 21,1 \% & 21,1 \% \\
\bottomrule
\end{tabular}
\end{center}
\begin{center}\scriptsize
\begin{tabular}{@{}l r r r r r@{}}
\toprule
notre réplique (2014-26) & excès/an & $t$ & $\alpha$/an & $\beta$ & NAV \\
\midrule
Congress Buys $\cdot\ \tau$ & $+0{,}84\,\%$ & 0,74 & $-0{,}21\,\%$ & 1,06 & 565 \\
Congress Buys $\cdot\ \delta$ & $-0{,}03\,\%$ & $-0{,}03$ & $-0{,}69\,\%$ & 1,04 & 511 \\
Long-Short $\cdot\ \tau$ & $+1{,}57\,\%$ & 1,13 & $+0{,}32\,\%$ & 1,07 & 606 \\
Long-Short $\cdot\ \delta$ & $+0{,}92\,\%$ & 0,60 & $+0{,}07\,\%$ & 1,04 & 555 \\
\bottomrule
\end{tabular}
\end{center}

\textbf{1bis. Quiver --- version \textit{House uniquement}} (page \og{}U.S. House Long-Short\fg{}, même mécanique, flux restreint à la Chambre). \textbf{Ce qu'ils annoncent} (depuis 2020-04) : \textbf{$+$628,58 \%} \og{}All Time\fg{} $\cdot$ CAGR \textbf{36,94 \%} $\cdot$ $\beta$ 1,09 $\cdot$ Sharpe 1,024 $\cdot$ IR 0,71 $\cdot$ MaxDD $-$24,70 \%. \textbf{Notre réplique fidèle} (House, $\delta$, base 100 au 2020-04) : \textbf{$+$267 \%} $\cdot$ CAGR \textbf{23,2 \%} $\cdot$ $\beta$ 1,04 $\cdot$ Sharpe 0,74 $\cdot$ IR 0,17 $\cdot$ MaxDD $-$31,5 \% \emph{(SPY même départ : $+$230 \%)}. Un contrôle d'identité vérifie qu'en retirant le filtre chambre on retombe exactement sur la réplique Congress ci-dessus.
\begin{center}\scriptsize
\begin{tabular}{@{}l r r r r r r@{}}
\toprule
réplique House (2014-26) & excès/an & $t$ & $\alpha$/an & $\beta$ & Sharpe & NAV \\
\midrule
Congress Buys $\cdot\ \delta$ & $+0{,}61\,\%$ & 0,48 & $-0{,}18\,\%$ & 1,04 & 0,72 & 546 \\
Long-Short $\cdot\ \tau$ & $+1{,}44\,\%$ & 1,09 & $+0{,}28\,\%$ & 1,07 & 0,74 & 600 \\
Long-Short $\cdot\ \delta$ & $+1{,}48\,\%$ & 0,86 & $+0{,}50\,\%$ & 1,04 & 0,74 & 587 \\
Long-Short $\cdot\ \delta\ $ hebdo & $+1{,}68\,\%$ & 1,04 & $+0{,}64\,\%$ & 1,05 & 0,75 & 609 \\
\bottomrule
\end{tabular}
\end{center}
\textbf{Même verdict que la version Congress} : la Chambre ne change rien --- $\approx$ le marché avec du levier, $\alpha\approx 0$. L'écart au 36,94 \% tient à leur boîte noire non divulguée (fenêtre, timing $\tau/\delta$, panier short) et à un panneau de métriques incohérent ($\beta$ 1,09 impose une vol $\ge$ 19 \%/an, incompatible avec leur vol affichée).

\fig{quiver.png}{}

\textbf{2. NANC / GOP} (ETF Subversive) --- \emph{trois lectures d'un même prospectus}
\begin{pts}
\item Leur prospectus est \textbf{qualitatif} : il dit \emph{quoi} surpondérer (gros achats, récurrents, de plusieurs élus), jamais \emph{combien}.
\item Le fonds a \textbf{changé de mécanique} en août 2024 (B1bis). On garde donc \textbf{trois} lectures : \textbf{A} le score (mensuel), \textbf{B} le livre (événementiel), et \textbf{C} $=$ \textbf{A puis B}.
\item Résultat : $\sim$50 \% d'exposition partagée avec le vrai NANC, $\sim$30 \% avec le vrai GOP. \textbf{C} est la référence, \textbf{aucune} n'est meilleure sur tous les critères. Le \textbf{B8} mesure pourquoi.
\end{pts}

\textbf{B1 $\cdot$ Le socle commun aux trois}
\begin{pts}
\item Univers : actions cotées tradées par les élus du parti $P$ (Dém.\ $\to$ NANC, Rép.\ $\to$ GOP) via les PTR.
\item Fenêtre \textbf{3 ans glissants sur la divulgation} $\delta$ (seule date réplicable) : $\mathcal{K}(t)=\{k:\ t-3\,\text{ans}<\delta_k\le t\}$.
\item Montant $A_k$ $=$ \textbf{milieu de la fourchette} PTR (imposé par le prospectus) ; exécution au \textbf{close J+1}, noté $e_k$.
\item Classes d'actions fusionnées ($\Phi$ : GOOGL$\to$GOOG, BRK-A$\to$BRK-B\ldots), des deux côtés de la comparaison.
\end{pts}
\begin{center}\scriptsize
\begin{tabular}{@{}l l@{}}
\toprule
symbole & définition \\
\midrule
$k$ $\cdot$ $m$ $\cdot$ $i$ & transaction $\cdot$ élu $\cdot$ titre \\
$\varepsilon_k=\pm1$ & achat $/$ vente \\
$B_i=\sum_{\text{achats}}A_k$ & \$ acheté sur la fenêtre \\
$m_i$ & nb d'élus acheteurs distincts \\
$\gamma_k\in[0,1]$ & part \emph{récente} d'une vente (B4) \\
$s_i$ $\cdot$ $w_i$ $\cdot$ $q_i$ & score $\cdot$ poids $\cdot$ nb de \textbf{parts} \\
$\Pi_c$ & plafonner à $c$, puis renormaliser \\
$N_{\mathrm{eff}}=1/\sum_i w_i^2$ & nb \emph{effectif} de lignes \\
$w^{\star}_i$ & poids \textbf{publié} par le vrai fonds \\
\bottomrule
\end{tabular}
\end{center}

\textbf{B1bis $\cdot$ Ce qu'ils publient vraiment} \emph{(prospectus statutaire $+$ SAI, rapports annuels N-CSR)}
\begin{center}\scriptsize
\begin{tabular}{@{}l r r@{}}
\toprule
 & NANC & GOP \\
\midrule
création & \multicolumn{2}{c}{06/02/2023} \\
frais $\cdot$ AUM & 0,72 \% $\cdot$ 282 M\$ & 0,73 \% $\cdot$ 88 M\$ \\
lignes & 104 & 138 \\
\addlinespace[1pt]
\emph{depuis création}, /an & $+23{,}63\,\%$ & $+14{,}75\,\%$ \\
S\&P 500 TR (réf.) & \multicolumn{2}{c}{$+20{,}96\,\%$} \\
\addlinespace[1pt]
\textbf{turnover} FY23 & 44 \% & 46 \% \\
\quad FY24 & 62 \% & 54 \% \\
\quad FY25 & \textbf{10 \%} & \textbf{16 \%} \\
\bottomrule
\end{tabular}
\end{center}
\begin{pts}
\item \textbf{Trois régimes, pas un} : \emph{Subversive Capital} (02/2023 $\to$ 07/2024, \textbf{500--600 lignes}), puis \textbf{Tidal} (dès le 02/08/2024, \textbf{100--200 lignes}), réorganisation 12/2024.
\item Le texte a changé : 2023 \og{}\emph{to create the Fund's \textbf{initial} portfolio}\fg{} $\to$ 2026 \og{}\emph{to \textbf{manage} the portfolio}\fg{}. Et la lettre du gérant FY2024 dit : \og{}\emph{\textbf{eliminated 582 very small positions}}\fg{} $\Rightarrow$ un \textbf{reset daté} au 02/08/2024.
\item \textbf{Data Provider $=$ Unusual Whales}, mais le prospectus précise qu'il \og{}\emph{plays no role in the construction}\fg{} : c'est \textbf{Tidal} qui construit --- et le rapport FY2025 avoue une surcouche \textbf{discrétionnaire} : \og{}\emph{cross referenced against the \ldots{} committee roles}\fg{}.
\item Preuves de \textbf{dérive} (N-CSR FY2025) : plus-value latente $+22{,}4\,\%$, gains/pertes latents \textbf{7,6:1}, AXP $+21{,}2\,\%$ de parts $=$ exactement l'effet des créations (\emph{zéro} arbitrage sur 12 mois). Mais NVDA $-15\,\%$ de parts $\Rightarrow$ un \textbf{plafonnement de fait}, nulle part écrit.
\item Ce qu'ils \textbf{n'utilisent pas} : les déclarations \emph{annuelles} de patrimoine (\og{}\emph{acquired prior to swearing in \ldots{} are not considered}\fg{}) --- vérifié chez nous : AMAT pèse \textbf{0,009 \%} des patrimoines démocrates contre 5,19 \% du fonds.
\end{pts}

\textbf{B2 $\cdot$ Deux mécaniques opposées}
\begin{center}\scriptsize
\begin{tabular}{@{}R{17mm} R{25mm} R{25mm}@{}}
\toprule
 & \textbf{A --- le score} & \textbf{B --- le livre} \\
\midrule
clause lue & ventes ignorées & on suit les divulgations \\
ventes & aucun effet & retirent des parts \\
objet & score $s_i$ & livre en parts $q_i$ \\
univers & top-150 & tout titre déclaré \\
plafond $c$ & 8,5 \% & \textbf{aucun} \\
recomposition & mensuelle & \textbf{aucune} \\
si $P_i$ monte & ramené à la cible & \textbf{poids augmente} \\
\midrule
turnover & 66 \% & 13 \% \\
\textbf{régime décrit} & \textbf{Subversive} 2023-24 (44--62 \%) & \textbf{Tidal} 2025 (10 \%) \\
\bottomrule
\end{tabular}
\end{center}
\lire{les deux colonnes sont les deux \emph{moitiés} de la vie du fonds, pas deux candidats concurrents. \textbf{C} les enchaîne.}

\textbf{B3 $\cdot$ A --- le score} \emph{(recomposé chaque mois)}
\begin{pts}
\item \textbf{Lecture} : le prospectus interdit d'ajuster sur les ventes de titres acquis \emph{avant} le suivi $\Rightarrow$ on les ignore \textbf{toutes}.
\item \textbf{Score} $=$ taille $\times$ largeur ; \textbf{sélection} $=$ les 150 plus gros ; plafond $c=8{,}5\,\%$ ; recomposé chaque mois.
\end{pts}
\[ s_i=B_i\,m_i \qquad w=\Pi_c\Bigl(s\,\big/\,\textstyle\sum_{j\in\mathcal{U}}s_j\Bigr) \]

\textbf{B4 $\cdot$ B --- le livre} \emph{(événementiel, il dérive)}
\begin{pts}
\item \textbf{Lecture} : \og{}\emph{the Fund will buy or sell a security \textbf{when} a position is reported}\fg{} --- pas un score recalculé, un \textbf{livre} alimenté par les divulgations.
\item \textbf{Livre initial} à la création $t_0=$ 06/02/2023 : le net des 3 ans précédents, converti en \textbf{parts} au prix de $t_0$.
\item \textbf{Ensuite} : un achat ajoute $A_k/P_i(e_k)$ parts, une vente en retire $\gamma_kA_k/P_i(e_k)$, où $\gamma_k$ est la part \emph{récente} de la vente (une vente qui solde un lot dormant a $\gamma_k\!\to\!0$ : elle \textbf{n'ajuste rien}, comme l'exige le prospectus).
\item \textbf{Ni recomposition, ni plafond, ni sélection} : entre deux divulgations les poids \textbf{dérivent avec les prix}.
\end{pts}
\[ q_i(t)=\frac{\mathrm{net}_i(t_0)}{P_i(t_0)}+\!\!\sum_{k\,:\,t_0<e_k\le t}\!\!\varepsilon_k\,\gamma_k\,\frac{A_k}{P_i(e_k)} \]
\[ w_i(t)=q_i(t)P_i(t)\,\big/\,\textstyle\sum_j q_j(t)P_j(t) \]

\textbf{B5 $\cdot$ C --- la bascule} \emph{(la référence)}
\begin{pts}
\item \textbf{A jusqu'au 02/08/2024}, puis le portefeuille est converti en \textbf{parts} et bascule en \textbf{B} : $q_i(t_1)=w^{A}_i\,S/P_i$, avec $S$ la taille du livre en \$.
\item C'est la seule lecture qui suive \textbf{l'histoire réelle} du fonds --- et la seule qui rende le \textbf{GOP} exploitable.
\end{pts}
\obs{\emph{Piège} : convertir des poids ($\sum w=1$) en parts \textbf{sans les mettre à l'échelle en dollars} donne un livre valant 1 \$ --- un seul achat de quelques milliers de \$ écrase alors tout. Il faut calibrer $S$ sur la taille du régime.}

\textbf{B5bis $\cdot$ Trois tests} \emph{(écartés comme versions, gardés comme enseignements)}
\begin{pts}
\item \textbf{Netter les ventes récentes} (purge $\gamma$ appliquée au \emph{score}) \textbf{ronge l'excès de NANC} : $+2{,}98\to+1{,}40\,\%$/an, $\alpha$ $+1{,}26\to-0{,}50$. $\Rightarrow$ le $+2{,}98$ de \textbf{A} est en partie l'artefact du \og{}on ignore les ventes\fg{}.
\item \textbf{Le plafond n'est utile qu'à A} : en mécanique \textbf{B}, \textbf{NVDA sort à 8,9 \% tout seul} (réel 8,5) ; en mécanique \textbf{A} sans plafond il monte à \textbf{37 \%} ($N_{\mathrm{eff}}$ 6 contre 33 réel). Le fonds lui-même écrête (NVDA $-15\,\%$ de parts, B1bis).
\item \textbf{Élaguer en continu} à 150 lignes règle le cardinal (719 $\to$ 150, réel 104) et la rotation (26 $\to$ 21 \%), \textbf{mais casse le classement} ($\rho$ de rang $+0{,}32\to-0{,}02$). $\Rightarrow$ \textbf{nos 150 plus grosses lignes ne sont pas leurs 104} : c'est notre \emph{ordre de grandeur relatif} qui est faux, pas la longueur de la liste.
\end{pts}

\textbf{B6 $\cdot$ Ce qu'on obtient} --- (a) fidélité aux holdings publiés :
\begin{center}\scriptsize
\begin{tabular}{@{}l c c r r r r@{}}
\toprule
 & top-10 & top-15 & $\rho_{\text{P}}$ & $\rho_{\text{S}}$ & expo. & recouv. \\
\midrule
NANC A & 5/10 & 7/15 & 0,73 & 0,06 & \textbf{51 \%} & \textbf{65 \%} \\
NANC B & 5/10 & 8/15 & 0,74 & 0,33 & 46 \% & 51 \% \\
NANC \textbf{C} & 5/10 & 7/15 & \textbf{0,75} & 0,32 & 48 \% & 57 \% \\
\addlinespace[1pt]
GOP A & 2/10 & 5/15 & 0,15 & $-0{,}15$ & \textbf{31 \%} & \textbf{25 \%} \\
GOP B & 0/10 & 0/15 & 0,16 & 0,36 & 26 \% & 14 \% \\
GOP \textbf{C} & 1/10 & 2/15 & \textbf{0,26} & 0,28 & 29 \% & 20 \% \\
\bottomrule
\end{tabular}
\end{center}
\lire{$\rho_{\text{P}}$ Pearson sur les poids, $\rho_{\text{S}}$ \textbf{sur les rangs} ; expo. $=\sum_i\min(w_i,w^{\star}_i)$ ; recouv. $=$ la même somme restreinte à leur top-10.}
\obs{$\rho_{\text{P}}\approx0{,}7$ est \textbf{porté par NVDA seul} : sur les \textbf{rangs}, la corrélation tombe à $\mathbf{0{,}06}$ pour A. On retrouve leurs \emph{gros noms}, pas leur \emph{classement}.}
\obs{\textbf{C} obtient le meilleur $\rho$ et \textbf{sauve le GOP} (0/10 $\to$ 1/10, $\rho$ 0,16 $\to$ 0,26) ; \textbf{A} garde le meilleur recouvrement du top-10 (65 \%) et la meilleure exposition (51 \%). \textbf{Aucune ne domine sur tous les critères} --- on l'assume.}

(a\,bis) \textbf{le turnover} --- test de la \emph{mécanique}, pas de la composition :
\begin{center}\scriptsize
\begin{tabular}{@{}l r r r r@{}}
\toprule
turnover/an & 2023 & 2024 & 2025 & médiane \\
\midrule
A (mensuel) & 59 \% & 67 \% & 82 \% & 66 \% \\
B (livre) & 13 \% & 11 \% & 16 \% & 13 \% \\
\textbf{C} (bascule) & 48 \% & 37 \% & 26 \% & --- \\
\addlinespace[1pt]
\textbf{fonds NANC} & \textbf{44 \%} & \textbf{62 \%} & \textbf{10 \%} & --- \\
\bottomrule
\end{tabular}
\end{center}
\lire{turnover $=\min(\text{achats},\text{ventes})/\text{actif}$ (définition SEC) $=\frac12\sum_i|w^{\text{cible}}_i-w^{\text{dérivé}}_i|$ cumulé sur l'année. Le chiffre du fonds exclut les échanges \emph{in kind} : c'est un plancher.}
\obs{\textbf{C'est la découverte du dossier.} Le fonds tournait à \textbf{44--62 \%} en 2023-24 $\approx$ notre \textbf{A} (59--67 \%) ; il est tombé à \textbf{10 \%} en 2025 $\approx$ notre \textbf{B} (11--16 \%). \emph{Aucune version unique ne peut coller aux deux} --- d'où \textbf{C}.}

(b) \textbf{composition} --- ancrée sur le top-10 du vrai NANC ($w^{\star}$), en \% :
\begin{center}\scriptsize
\begin{tabular}{@{}l r r r r r@{}}
\toprule
NANC & $w^{\star}$ & A & B & C & $\Delta$ \\
\midrule
NVDA & 8,53 & 8,50 & 8,93 & 11,73 & 0,03 \\
GOOG & 5,81 & 8,50 & 4,96 & 8,94 & 0,85 \\
MSFT & 5,29 & 8,22 & 3,81 & 5,55 & 0,26 \\
\textbf{AMAT} & 5,19 & 0,21 & 0,31 & 0,42 & \textbf{4,77} \\
AMZN & 4,70 & 8,24 & 3,05 & 2,88 & 1,65 \\
AAPL & 4,58 & 6,57 & 0,61 & 2,38 & 1,99 \\
\textbf{CRWD} & 3,41 & 0,13 & 1,38 & 0,07 & \textbf{2,03} \\
\textbf{PM} & 3,20 & 0,11 & 0,05 & 0,04 & \textbf{3,09} \\
AXP & 2,84 & 0,15 & 0,65 & 0,79 & 2,05 \\
LLY & 2,71 & 0,75 & 0,43 & 0,36 & 1,96 \\
\midrule
\multicolumn{6}{@{}l@{}}{\footnotesize $\Sigma$ top-10 réel $=46{,}3$ \% $\cdot$ \textbf{recouvrement} $\sum_i\min(w,w^{\star})$ :} \\
\multicolumn{6}{@{}l@{}}{\footnotesize \quad A \textbf{65 \%} $\cdot$ B \textbf{51 \%} $\cdot$ C \textbf{57 \%} du bloc} \\
\bottomrule
\end{tabular}
\end{center}
\lire{$\Delta=\min_v|w_v-w^{\star}|$ : l'écart de la version la \emph{plus} proche --- donc une borne basse de notre erreur sur cette ligne. Le recouvrement est \textbf{non compensatoire} : surpondérer GOOG ne rachète pas AMAT manquant.}

Et l'erreur \textbf{symétrique} --- nos gros poids qu'ils n'ont \emph{pas} :
\begin{center}\scriptsize
\begin{tabular}{@{}l r r r r@{}}
\toprule
 & $w^{\star}$ & A & B & C \\
\midrule
MU & 2,60 & 0,62 & \textbf{7,73} & 3,40 \\
DIS & 1,04 & 0,47 & 1,57 & \textbf{4,55} \\
\bottomrule
\end{tabular}
\end{center}

Et le \textbf{GOP}, où rien ne colle :
\begin{center}\scriptsize
\begin{tabular}{@{}l r r r r r@{}}
\toprule
GOP & $w^{\star}$ & A & B & C & $\Delta$ \\
\midrule
\textbf{FIX} & 8,44 & 0,00 & 0,52 & 0,01 & \textbf{7,92} \\
INTC & 5,51 & 1,67 & 1,04 & 2,90 & 2,61 \\
JPM & 4,47 & 1,34 & 0,94 & 1,50 & 2,97 \\
NVDA & 3,22 & 6,42 & 0,93 & 1,70 & 1,52 \\
ANET & 3,22 & 0,67 & 0,36 & 0,01 & 2,55 \\
IBIT & 3,00 & \multicolumn{4}{c}{\emph{hors univers} (ETF crypto)} \\
AMD & 2,95 & 2,17 & 0,55 & 1,03 & 0,78 \\
UTHR & 2,07 & 0,00 & 0,43 & 0,00 & 1,64 \\
\midrule
\multicolumn{6}{@{}l@{}}{\footnotesize recouvrement : A \textbf{25 \%} $\cdot$ B \textbf{14 \%} $\cdot$ C \textbf{20 \%} du bloc} \\
\bottomrule
\end{tabular}
\end{center}
\fig{nanc_gop_v3.png}{réel (gris) vs mécanique \textbf{B} : \textbf{NVDA émerge à 8,9 \%} sans aucun plafond (réel 8,5 --- la mécanique s'auto-régule) ; AMAT reste à 0,3 (réel 5,2) et FIX à 0,5 (réel 8,4).}

(c) \textbf{concentration} ($N_{\mathrm{eff}}=1/\sum w_i^2$) :
\begin{center}\scriptsize
\begin{tabular}{@{}l r r r r@{}}
\toprule
 & max $w$ & top-5 & $N_{\mathrm{eff}}$ & titres \\
\midrule
NANC \textbf{réel} & 8,5 & 29,5 & \textbf{33} & \textbf{104} \\
NANC A & 8,5 & 40,0 & 23 & 150 \\
NANC B & 8,9 & 28,5 & 45 & 1\,239 \\
NANC \textbf{C} & 11,7 & 35,1 & \textbf{29} & 719 \\
\bottomrule
\end{tabular}
\end{center}
\obs{\textbf{C} a le $N_{\mathrm{eff}}$ le plus proche du réel (29 contre 33) --- mais avec 719 lignes au lieu de 104 : sa concentration est juste \emph{par hasard}, en étalant beaucoup de très petites lignes. Le $N_{\mathrm{eff}}$ de A, lui, est fixé par le \textbf{plafond}, pas par le signal.}

(d) \textbf{performance} (excès $=$ rendement $-$ SPY) :
\begin{center}\scriptsize
\begin{tabular}{@{}l l r r r r@{}}
\toprule
 & période & exc./an & $\alpha$/an & $\beta$ & NAV \\
\midrule
NANC A & 2014-26 & $+2{,}98$ & $+1{,}26$ & 1,08 & 672 \\
GOP A & 2014-26 & $+1{,}24$ & $+1{,}12$ & 1,00 & 612 \\
\addlinespace[1pt]
NANC B & 2023-26 & $-0{,}21$ & $-1{,}52$ & 1,03 & 197 \\
GOP B & 2023-26 & $-1{,}47$ & $-0{,}14$ & 0,89 & 189 \\
NANC \textbf{C} & 2023-26 & $\mathbf{+2{,}46}$ & $+1{,}46$ & 1,06 & 213 \\
GOP \textbf{C} & 2023-26 & $-0{,}27$ & $+1{,}13$ & 0,93 & 196 \\
\bottomrule
\end{tabular}
\end{center}
\lire{NAV base 100 au \textbf{début de la période} de la ligne --- 197 (3 ans) et 672 (12 ans) ne se comparent pas. SPY $=517$ sur 2014-26.}
\obs{Sur sa propre période, \textbf{C} retrouve l'ordre de grandeur du vrai NANC : $+2{,}46\,\%$/an d'excès contre $+2{,}67$ réel. Mais $\beta\approx1$ et $t_\alpha<1$ partout : c'est du \textbf{marché}, pas de l'alpha.}
\fig{nanc_gop.png}{}
\fig{nanc_rolling.png}{$\beta$ et $\alpha$ glissants (OLS 252 j, validé $=$ \texttt{polyfit}) : $\beta$ colle à $\sim$1 ; l'$\alpha$ traverse zéro sans cesse.}

\textbf{B7 $\cdot$ Ce qui n'explique \emph{pas} l'écart} --- on a audité et réparé notre propre chaîne :
\begin{pts}
\item \textbf{7 287 transactions OCR réintégrées} (892 PTR papier écartés par le filtre anti-manuscrit, dont 6 482 républicaines --- McCaul 2015-16).
\item \textbf{10 séries de prix} récupérées (renommages MMC$\to$MRSH, FI$\to$FISV, NLOK$\to$GEN\ldots).
\item Montants ouverts (\og{}over \$1M\fg{}, 412 lignes au plancher) revalorisés à 3 M\$.
\item $\Rightarrow$ \textbf{la fidélité ne bouge pas} ($\rho$ identiques à $0{,}01$ près) ; la revalorisation \emph{dégrade} même GOP. \textbf{Aucun} des PTR récupérés ne contient un flux à l'échelle du FIX 8,4 \%.
\end{pts}

\textbf{B8 $\cdot$ Diagnostic --- leurs poids ne suivent pas le flux déclaré}
\[ \lambda_i \;=\; w^{\star}_i \;\Big/\; \Bigl(F_i\big/\textstyle\sum_j F_j\Bigr),\qquad F_i=\text{\$ d'achats déclaré sur }i \]
\begin{center}\scriptsize
\begin{tabular}{@{}l r@{\ }r@{\ }r@{\ }r@{\ }r@{\ }r@{\ }r@{}}
\toprule
NANC & MSFT & AMZN & NVDA & AXP & AMAT & CRWD & PM \\
$\lambda_i$ & 0,5 & 2,4 & 2,7 & 9,2 & 17,0 & 18,1 & \textbf{22,9} \\
\bottomrule
\end{tabular}
\end{center}
\begin{pts}
\item \textbf{Ce qu'on mesure} : si le fonds pondérait proportionnellement au flux qu'on observe, $\lambda_i$ serait \textbf{constant}.
\item \textbf{Ce qu'on observe} : $\lambda$ va de 0,5 à 22,9 --- amplitude $\times$\textbf{50} (GOP : $\times$29, avec FIX à 247).
\item \textbf{Aucun observable ne les ordonne} : corrélations \emph{de rang} de $w^{\star}$ avec le \$ \textbf{0,19}, le nb de trades \textbf{0,15}, le nb d'élus \textbf{0,18}.
\item \textbf{Portée} : nos paramètres ($K$, $\theta$, $c$, fenêtre) agissent \textbf{uniformément} sur le vecteur de poids --- ils ne peuvent pas créer une dispersion $\times$50 titre par titre. Avec B7, l'écart n'est imputable ni à nos données ni à notre paramétrage : il porte sur une information \textbf{absente des PTR} (leur \emph{Data Provider}, et la \textbf{discrétion} que le prospectus accorde explicitement à l'adviser).
\end{pts}

\end{multicols}

% ═══════════ ANNEXE : extraits sources ═══════════
\vspace{2mm}\hrule\vspace{1.5mm}
{\color{bleu}\bfseries\large Annexe --- extraits sources}\par\vspace{2mm}

\textbf{1. Quiver --- Congress Buys} \emph{(quiverquant.com)} :
{\footnotesize\itshape ``The Congress Buys Strategy tracks the performance of stocks that have been purchased by members of U.S. Congress (or their family). This strategy is weighted based on the reported size of the purchases, with weekly rebalancing.''\par}
\vspace{1.5mm}

\textbf{1. Quiver --- Congress Long-Short} \emph{(quiverquant.com)} :
{\footnotesize\itshape ``The Congress Long-Short Strategy tracks the performance of stocks that have been purchased or sold by members of U.S. Congress (or their family). It takes a long position in stocks that have been purchased, and a short position in stocks which have been sold. This strategy is weighted based on the reported size of the transactions and employs leverage with 130\% long exposure and 30\% short exposure, with weekly rebalancing.''\par}
\vspace{2mm}

\textbf{2. NANC} --- \emph{prospectus \og{}Unusual Whales Subversive Democratic Trading ETF\fg{}, \og{}Principal Investment Strategies\fg{}, p.\,2} :
{\footnotesize\itshape
``To manage the Fund's portfolio, the Adviser will obtain and use information derived by the Data Provider from PTRs filed by Democratic U.S. Congresspeople for the past 3 years. Purchases made during that time will be netted against any sales of the same security when managing the Fund's portfolio of equity securities. As the investment thesis of the Fund is to track the trading activity of Democratic U.S. Congresspeople while in office, equity securities acquired by Democratic U.S. Congresspeople prior to his or her swearing in (or the 3-year lookback period) are not considered when managing the Fund's portfolio. To the extent a Democratic U.S. Congressperson sells equity securities that were acquired prior to his or her swearing in, the Adviser will not adjust the Fund's portfolio.

Under normal circumstances, the Fund will invest in a portfolio of between 100 to 200 holdings. However, the number and size of positions held by the Fund will vary based on the number of positions traded by Democratic U.S. Congresspeople. Individual holdings will be weighted based on the level of reported trading in a security by Democratic U.S. Congresspeople. Securities with large purchases, recurring purchases and purchases from multiple Democratic U.S. Congresspeople will be overweighted. Securities with small purchases, recent sales and one-off trades by Democratic U.S. Congresspeople will be excluded or underweighted. The Adviser may exclude positions that are traded by Democratic U.S. Congresspeople at its discretion. Considerations for excluding a security include, but are not limited to, limited liquidity of such security and pending corporate actions that may impact such security.''\par}
\vspace{2mm}

\textbf{2. GOP} --- \emph{prospectus \og{}Unusual Whales Subversive Republican Trading ETF\fg{}, \og{}Principal Investment Strategies\fg{}, p.\,2} :
{\footnotesize\itshape
``To manage the Fund's portfolio, the Adviser will obtain and use information derived by the Data Provider from PTRs filed by Republican U.S. Congresspeople for the past 3 years. Purchases made during that time will be netted against any sales of the same security when managing the Fund's portfolio of equity securities. As the investment thesis of the Fund is to track the trading activity of Republican U.S. Congresspeople while in office, equity securities acquired by Republican U.S. Congresspeople prior to his or her swearing in (or the 3-year lookback period) are not considered when managing the Fund's portfolio. To the extent a Republican U.S. Congressperson sells equity securities that were acquired prior to his or her swearing in, the Adviser will not adjust the Fund's portfolio.

Under normal circumstances, the Fund will invest in a portfolio of between 100 to 200 holdings. However, the number and size of positions held by the Fund will vary based on the number of positions traded by Republican U.S. Congresspeople. Individual holdings will be weighted based on the level of reported trading in a security by Republican U.S. Congresspeople. Securities with large purchases, recurring purchases and purchases from multiple Republican U.S. Congresspeople will be overweighted. Securities with small purchases, recent sales and one-off trades by Republican U.S. Congresspeople will be excluded or underweighted. The Adviser may exclude positions that are traded by Republican U.S. Congresspeople at its discretion. Considerations for excluding a security include, but are not limited to, limited liquidity of such security and pending corporate actions that may impact such security.''\par}

\end{document}
